\documentclass[10pt,a4paper]{amsart}
\usepackage{amsmath,amssymb,amsthm,graphicx}
% Retained preamble for the Stacks Project chapter snapshot in this directory.
% The chapter's own preamble.tex is a copy of the shared Stacks Project preamble; the block
% below is sliced MECHANICALLY, line for line, from it. Nothing was added, removed or
% redefined; the wrapper class, the xr-hyper/hyperref wiring and the externaldocument list
% are omitted because the assembler supplies the document class.
%
% lines taken from preamble.tex: 3, 5, 9, 11, 12, 156, 157, 158, ... (51 lines total)
%
% package coverage restored from the chapter's own preamble (2026-10-13): xy, xy-2cell

\usepackage[all]{xy}
\xyoption{2cell}
\UseAllTwocells

\newenvironment{reference}{\comment}{\endcomment}
\newenvironment{slogan}{\comment}{\endcomment}
\newenvironment{history}{\comment}{\endcomment}
\theoremstyle{plain}
\newtheorem{theorem}[subsection]{Theorem}
\newtheorem{proposition}[subsection]{Proposition}
\newtheorem{lemma}[subsection]{Lemma}
\theoremstyle{definition}
\newtheorem{definition}[subsection]{Definition}
\newtheorem{example}[subsection]{Example}
\newtheorem{exercise}[subsection]{Exercise}
\newtheorem{situation}[subsection]{Situation}
\theoremstyle{remark}
\newtheorem{remark}[subsection]{Remark}
\newtheorem{remarks}[subsection]{Remarks}
\numberwithin{equation}{subsection}
\def\lim{\mathop{\mathrm{lim}}\nolimits}
\def\colim{\mathop{\mathrm{colim}}\nolimits}
\def\Spec{\mathop{\mathrm{Spec}}}
\def\Hom{\mathop{\mathrm{Hom}}\nolimits}
\def\Ext{\mathop{\mathrm{Ext}}\nolimits}
\def\SheafHom{\mathop{\mathcal{H}\!\mathit{om}}\nolimits}
\def\SheafExt{\mathop{\mathcal{E}\!\mathit{xt}}\nolimits}
\def\Sch{\mathit{Sch}}
\def\Mor{\mathop{\mathrm{Mor}}\nolimits}
\def\Ob{\mathop{\mathrm{Ob}}\nolimits}
\def\Sh{\mathop{\mathit{Sh}}\nolimits}
\def\NL{\mathop{N\!L}\nolimits}
\def\CH{\mathop{\mathrm{CH}}\nolimits}
\def\proetale{{pro\text{-}\acute{e}tale}}
\def\etale{{\acute{e}tale}}
\def\QCoh{\mathit{QCoh}}
\def\Ker{\mathop{\mathrm{Ker}}}
\def\Im{\mathop{\mathrm{Im}}}
\def\Coker{\mathop{\mathrm{Coker}}}
\def\Coim{\mathop{\mathrm{Coim}}}
\DeclareMathSymbol{\boxtimes}{\mathbin}{AMSa}{"02}
\def\QCohstack{\mathcal{QC}\!\mathit{oh}}
\def\Cohstack{\mathcal{C}\!\mathit{oh}}
\def\Spacesstack{\mathcal{S}\!\mathit{paces}}
\def\Quotfunctor{\mathrm{Quot}}
\def\Hilbfunctor{\mathrm{Hilb}}
\def\Curvesstack{\mathcal{C}\!\mathit{urves}}
\def\Polarizedstack{\mathcal{P}\!\mathit{olarized}}
\def\Complexesstack{\mathcal{C}\!\mathit{omplexes}}
\def\Pic{\mathop{\mathrm{Pic}}\nolimits}
\def\Picardstack{\mathcal{P}\!\mathit{ic}}
\def\Picardfunctor{\mathrm{Pic}}
\def\Deformationcategory{\mathcal{D}\!\mathit{ef}}
\usepackage{cleveref}
\crefname{originallemma}{original lemma}{original lemmas}
\Crefname{originallemma}{Original lemma}{Original lemmas}
\makeatletter
\newlabel{criteria-lemma-check-property-limit-preserving}{{07WI}{source}{}{sourcepointer-1}{}}
\newlabel{criteria-lemma-check-property-limit-preserving@cref}{{[originallemma][0][]07WI}{[1][0][]source}{}{}{}}
\makeatother
\crefname{originallemma}{original lemma}{original lemmas}
\Crefname{originallemma}{Original lemma}{Original lemmas}
\makeatletter
\newlabel{criteria-lemma-stacks-etale}{{076V}{source}{}{sourcepointer-2}{}}
\newlabel{criteria-lemma-stacks-etale@cref}{{[originallemma][0][]076V}{[1][0][]source}{}{}{}}
\makeatother
\crefname{originallemma}{original lemma}{original lemmas}
\Crefname{originallemma}{Original lemma}{Original lemmas}
\makeatletter
\newlabel{lemma-get-smooth}{{07XY}{source}{}{sourcepointer-3}{}}
\newlabel{lemma-get-smooth@cref}{{[originallemma][0][]07XY}{[1][0][]source}{}{}{}}
\makeatother
\crefname{originallemma}{original lemma}{original lemmas}
\Crefname{originallemma}{Original lemma}{Original lemmas}
\makeatletter
\newlabel{morphisms-lemma-enough-finite-type-points}{{02J4}{source}{}{sourcepointer-4}{}}
\newlabel{morphisms-lemma-enough-finite-type-points@cref}{{[originallemma][0][]02J4}{[1][0][]source}{}{}{}}
\makeatother
\crefname{originallemma}{original lemma}{original lemmas}
\Crefname{originallemma}{Original lemma}{Original lemmas}
\makeatletter
\newlabel{section-axioms}{{07XJ}{source}{}{sourcepointer-5}{}}
\newlabel{section-axioms@cref}{{[originallemma][0][]07XJ}{[1][0][]source}{}{}{}}
\makeatother
\crefname{originallemma}{original lemma}{original lemmas}
\Crefname{originallemma}{Original lemma}{Original lemmas}
\makeatletter
\newlabel{sets-lemma-what-is-in-it}{{000R}{source}{}{sourcepointer-6}{}}
\newlabel{sets-lemma-what-is-in-it@cref}{{[originallemma][0][]000R}{[1][0][]source}{}{}{}}
\makeatother
\title{The axioms imply X is an algebraic stack}
\author{The Stacks Project Authors}
\date{stacks-project:artin.tex:local-snapshot-20261012; GNU FDL 1.2 or later}
\begin{document}
\maketitle
\noindent\textit{Editorial note.} Statement span is lines 2728-2743 of sources/stacks-artin/pinned-original-source.tex; proof span is lines 2745-2780, the complete author proof (34 body lines).\par\medskip
\noindent\textit{Editorial note.} Navigation stubs: 6 cited-but-not-included labels are supplied as non-load-bearing pointers with their pinned tags where the snapshot tag index has them.\par\medskip
\noindent\textit{Editorial note.} extra\_packages is empty on purpose: the retained preamble already loads xy with options and a second declaration clashes in the assembler wrapper.\par\medskip
\noindent\textit{Editorial note.} Text is retained verbatim, including the author's own typos and placeholders.\par\medskip
\noindent\textit{Editorial note.} Queue index 22 of queue\_artin.json (pb=1877, stmt 2728-2743, proof 2745-2780).\par\medskip
\noindent\textit{Editorial note.} stacks-artin wave ar01 (2026-10-14): fresh adjudication accept per lane-s3/artin\_judgments.md; pool per Chief 10:5x direct dispatch (idle-first over S12's 01:5x note, SMC precedent); artin.tex 99 rows/67 with proof judged 46 accept / 21 reject / 5 AGP-overlap avoidance rows; patch-sd-envs contract spans scanned.\par\medskip
\noindent\textit{Allowed standard prerequisites (editorial).}\begin{itemize}
\item Explicit prerequisite boundary. The cited inputs are prerequisites; the argument in the retained span is not replaced by any cited result.
\end{itemize}\medskip
\ifdefined\hypertarget\hypertarget{sourcepointer-1}{}\fi
\noindent\textit{Source reference.} Original lemma 07WI: sources/stacks-artin/pinned-original-source.tex, source lines 2780--2780. Cited by the retained statement or proof but not part of the collected spans; supplied only as a navigation pointer. The argument does not depend on its content.\par\medskip
\ifdefined\hypertarget\hypertarget{sourcepointer-2}{}\fi
\noindent\textit{Source reference.} Original lemma 076V: sources/stacks-artin/pinned-original-source.tex, source lines 2780--2780. Cited by the retained statement or proof but not part of the collected spans; supplied only as a navigation pointer. The argument does not depend on its content.\par\medskip
\ifdefined\hypertarget\hypertarget{sourcepointer-3}{}\fi
\noindent\textit{Source reference.} Original lemma 07XY: sources/stacks-artin/pinned-original-source.tex, source lines 2408--2421. Cited by the retained statement or proof but not part of the collected spans; supplied only as a navigation pointer. The argument does not depend on its content.\par\medskip
\ifdefined\hypertarget\hypertarget{sourcepointer-4}{}\fi
\noindent\textit{Source reference.} Original lemma 02J4: sources/stacks-artin/pinned-original-source.tex, source lines 2780--2780. Cited by the retained statement or proof but not part of the collected spans; supplied only as a navigation pointer. The argument does not depend on its content.\par\medskip
\ifdefined\hypertarget\hypertarget{sourcepointer-5}{}\fi
\noindent\textit{Source reference.} Original lemma 07XJ: sources/stacks-artin/pinned-original-source.tex, source lines 2462--2489. Cited by the retained statement or proof but not part of the collected spans; supplied only as a navigation pointer. The argument does not depend on its content.\par\medskip
\ifdefined\hypertarget\hypertarget{sourcepointer-6}{}\fi
\noindent\textit{Source reference.} Original lemma 000R: sources/stacks-artin/pinned-original-source.tex, source lines 2780--2780. Cited by the retained statement or proof but not part of the collected spans; supplied only as a navigation pointer. The argument does not depend on its content.\par\medskip
\begin{lemma}
\label{lemma-diagonal-representable}
Let $S$ be a locally Noetherian scheme. Let
$p : \mathcal{X} \to (\Sch/S)_{fppf}$ be a category fibred in groupoids.
Assume that
\begin{enumerate}
\item $\Delta : \mathcal{X} \to \mathcal{X} \times \mathcal{X}$
is representable by algebraic spaces,
\item $\mathcal{X}$ satisfies axioms [-1], [0], [1], [2], [3] (see
Section \ref{section-axioms}),
\item every formal object of $\mathcal{X}$ is effective,
\item $\mathcal{X}$ satisfies openness of versality, and
\item $\mathcal{O}_{S, s}$ is a G-ring for all finite type points $s$ of $S$.
\end{enumerate}
Then $\mathcal{X}$ is an algebraic stack.
\end{lemma}

\begin{proof}
Lemma \ref{lemma-get-smooth} applies to $\mathcal{X}$. Using this we
choose, for every finite type field $k$ over $S$ and every
isomorphism class of object $x_0 \in \Ob(\mathcal{X}_{\Spec(k)})$,
an affine scheme $U_{k, x_0}$ of finite type over $S$ and a smooth morphism
$(\Sch/U_{k, x_0})_{fppf} \to \mathcal{X}$ such that there exists a finite
type point $u_{k, x_0} \in U_{k, x_0}$ with residue field $k$ such that $x_0$
is the image of $u_{k, x_0}$. Then
$$
(\Sch/U)_{fppf} \to \mathcal{X},
\quad\text{with}\quad
U = \coprod\nolimits_{k, x_0} U_{k, x_0}
$$
is smooth\footnote{Set theoretical remark: This coproduct is (isomorphic to)
an object of $(\Sch/S)_{fppf}$ as we have a bound on the index set
by axiom [-1], see Sets, Lemma \ref{sets-lemma-what-is-in-it}.}.
To finish the proof it suffices to show this map is surjective,
see Criteria for Representability, Lemma \ref{criteria-lemma-stacks-etale}
(this is where we use axiom [0]). By Criteria for Representability, Lemma
\ref{criteria-lemma-check-property-limit-preserving}
it suffices to show that
$(\Sch/U)_{fppf} \times_\mathcal{X} (\Sch/V)_{fppf} \to (\Sch/V)_{fppf}$
is surjective for those $y : (\Sch/V)_{fppf} \to \mathcal{X}$ where
$V$ is an affine scheme locally of finite presentation
over $S$. By assumption (1) the fibre product
$(\Sch/U)_{fppf} \times_\mathcal{X} (\Sch/V)_{fppf}$ is representable
by an algebraic space $W$. Then $W \to V$ is smooth, hence the image is
open. Hence it suffices to show that the image of $W \to V$ contains all
finite type points of $V$, see
Morphisms, Lemma \ref{morphisms-lemma-enough-finite-type-points}.
Let $v_0 \in V$ be a finite type point. Then $k = \kappa(v_0)$ is
a finite type field over $S$. Denote $x_0 = y|_{\Spec(k)}$
the pullback of $y$ by $v_0$. Then $(u_{k, x_0}, v_0)$ will give
a morphism $\Spec(k) \to W$ whose composition with $W \to V$
is $v_0$ and we win.
\end{proof}

\end{document}
